\section{The price of advice}\label{sec:advice}

The parent paper's second step from the one-hop model to machines is depth reduction
\cite[Section~7.4]{SeventeenthPaper2026}: a genie reveals values at chosen steps, and the revealed word is advice that
depends on position, so it breaks uniformity wherever it is used. The parent paper suggests that a version of
Theorem~TP with advice would charge each revealed bit as a regime change, at half a row per change, and proves nothing
about it. This section prices advice in the one-hop model: at most one dimension per symbol, and exactly that for
suitable advice, when it enters linearly; and over a finite field at most one dimension per symbol plus half a row for
each class change it induces.

\begin{definition}[Linear advice]\label{def:linear-advice}
\formal{Advice.lean}{IsClassUniformCodeWithAdvice}
Let $\alpha\colon\Mat_d(F)\to W$ be a linear map to any vector space. A subspace $\X$ is a \emph{code of a
class-uniform rule with linear advice $\alpha$} if there are decoders $\cA_J$, one for each class, and linear maps
$B_c\colon W\to F^{\Z/d}$, one for each row, such that $x_c=\cA_J(\sigma^cX)+B_c(\alpha(X))$ for every $X\in\X$, every
class $J$ and every $c\in J$.
\end{definition}

The maps $B_c$ may differ from row to row, so each row may use the advice in its own way.

\begin{theorem}[Theorem~C, linear advice costs its rank]\label{thm:linear-advice}
\formal{Advice.lean}{IsClassUniformCodeWithAdvice.inf_ker,finrank_le_finrank_inf_ker_add,IsClassUniformCodeWithAdvice.two_mul_finrank_le}
Let $\X$ be a code of a class-uniform rule with linear advice $\alpha$. Then $\X\cap\ker\alpha$ is a code of a
class-uniform rule for the same colouring, and $\dim\X\le\dim(\X\cap\ker\alpha)+\rank\alpha$. Consequently, if $S$ is as in
\cref{thm:TP}(b), then $2\dim\X\le d^2+d\card S+2\rank\alpha$; for $M\ge2$ interval classes,
$\dim\X\le(d^2+Md)/2+\rank\alpha$.
\end{theorem}

\begin{outline}
On $\X\cap\ker\alpha$ the advice term vanishes, so the same decoders work (\lean{IsClassUniformCodeWithAdvice.inf_ker}).
By rank and nullity for the restriction of $\alpha$ to $\X$, with values in the range of $\alpha$, the kernel of the
restriction, which is $\X\cap\ker\alpha$, has codimension at most $\rank\alpha$ in $\X$
(\lean{finrank_le_finrank_inf_ker_add}). The last statement adds \cref{thm:TP}
(\lean{IsClassUniformCodeWithAdvice.two_mul_finrank_le}).
\end{outline}

\begin{proposition}[The price is exact]\label{prop:sharp}
\formal{Advice.lean}{isClassUniformCodeWithAdvice_sup_range,finrank_sup_range}
Let $\X_0$ be a code of a class-uniform rule, and $p$ a linear projection of $\Mat_d(F)$ ($p^2=p$) that vanishes on
$\X_0$. Then $\X_0+\im p$ is a code of a class-uniform rule for the same colouring, with linear advice $p$, and
$\dim(\X_0+\im p)=\dim\X_0+\rank p$.
\end{proposition}

\begin{outline}
Take $B_c(w)=w_c-\cA_J(\sigma^cw)$ for $c\in J$. For $X=x+p(z)$ with $x\in\X_0$ we have $p(X)=p(z)$, and
$\cA_J(\sigma^cX)+B_c(p(X))=x_c+p(z)_c=X_c$ (\lean{isClassUniformCodeWithAdvice_sup_range}). The sum is direct, because
$y\in\X_0\cap\im p$ gives $y=p(y)=0$ (\lean{finrank_sup_range}).
\end{outline}

A projection of any rank $q\le d^2-\dim\X_0$ vanishing on $\X_0$ exists (project onto a $q$-dimensional subspace that
meets $\X_0$ only in $0$, along a complement containing $\X_0$). So every code of every class-uniform rule extends, with
linear advice of rank $q$, by exactly $q$ dimensions, for every $q\le d^2-\dim\X_0$: the inequality
$\dim\X\le\dim(\X\cap\ker\alpha)+\rank\alpha$ of \cref{thm:linear-advice} is sharp for every code. Its form with
\cref{thm:TP} is attained only from codes that attain \cref{thm:TP}.

Advice that is not linear may also change the decoders of the rows that use it. Over a finite field we can price that
as well.

\begin{definition}[Advice of any kind]\label{def:any-advice}
\formal{Advice.lean}{IsCodeWithAnyAdvice,ruleSpace}
Let $\Omega$ be a finite set, $w\colon\Mat_d(F)\to\Omega$ any map, and $\mathcal P_\omega$, for $\omega\in\Omega$,
colourings of $\Z/d$. A subspace $\X$ is a \emph{code with advice $w$} for the colourings $\mathcal P_\omega$ if for every
$\omega$ there are decoders $\cA^\omega_J$, one for each class $J$ of $\mathcal P_\omega$, and offsets
$b^\omega_c\in F^{\Z/d}$, one for each row, such that $x_c=\cA^\omega_J(\sigma^cX)+b^\omega_c$ for every $X\in\X$ with
$w(X)=\omega$, every class $J$ of $\mathcal P_\omega$ and every $c\in J$.
\end{definition}

The offsets let the advice enter every row in its own way, and the colouring $\mathcal P_\omega$ lets it change the
decoders of the rows that use it. A code of a class-uniform rule with linear advice $\alpha$ of rank $q$, over a finite
field, is a code with advice $w=\alpha$, $\abs{\Omega}=\abs F^q$ after restricting the values to the range of $\alpha$,
and $\mathcal P_\omega=\mathcal P$ for all~$\omega$.

\begin{theorem}[Theorem~D, the price of advice]\label{thm:any-advice}
\formal{Advice.lean}{IsCodeWithAnyAdvice.natCard_pow_finrank_le,IsCodeWithAnyAdvice.natCard_pow_two_mul_finrank_le,natCard_submodule}
Let $F$ be a finite field and $\X$ a code with advice $w\colon\Mat_d(F)\to\Omega$ for colourings $\mathcal P_\omega$.
\begin{enumerate}[label=(\roman*),leftmargin=2em]
\item If every code of a class-uniform rule for every $\mathcal P_\omega$ has dimension at most $D$, then
$\abs F^{\dim\X}\le\abs\Omega\cdot\abs F^{D}$.
\item If for every $\omega$ some non-empty set of at most $s$ rows contains every row whose class under
$\mathcal P_\omega$ differs from that of the row before it, then $\abs F^{2\dim\X}\le\abs\Omega^2\cdot\abs F^{d^2+ds}$.
\end{enumerate}
In particular, with $\abs\Omega=\abs F^q$, that is, $q$ symbols of advice, $2\dim\X\le d^2+ds+2q$.
\end{theorem}

\begin{outline}
Fix $\omega$ and a matrix $X_\omega$ of $\X$ with $w(X_\omega)=\omega$, if there is one. For every other such $X$, the
difference $X-X_\omega$ lies in $\X$ and satisfies the rule of $\omega$ without offsets, so it lies in the code
$\X\cap\{Y:y_c=\cA^\omega_J(\sigma^cY)\}$ (\lean{ruleSpace}), which has at most $\abs F^D$ elements
(\lean{natCard_submodule}). The fibres of $w$ partition $\X$, so $\abs\X\le\abs\Omega\cdot\abs F^D$
(\lean{IsCodeWithAnyAdvice.natCard_pow_finrank_le}). For (ii), \cref{thm:TP}(b) gives $D=\lfloor(d^2+ds)/2\rfloor$;
square (\lean{IsCodeWithAnyAdvice.natCard_pow_two_mul_finrank_le}).
\end{outline}

\begin{remark}[What this says about depth reduction]\label{rem:depth}
\Cref{thm:any-advice} is a version of \cref{thm:TP} with advice, in the one-hop model over a finite field. Every symbol
of advice costs one dimension, and the advice costs more only through the class changes it induces, at half a row
each. The parent paper suggests charging each revealed bit as a regime change, at half a row. Advice that enters
linearly, as it does when the revealed values are linear in the input and enter the rows additively, induces no class
change, and its price is one dimension per symbol (\cref{thm:linear-advice,prop:sharp}). Advice that changes the
decoder of a single row inside an interval induces two class changes, a full row. In the regime of depth reduction,
$q\approx\eps d^2$ symbols, part (ii) is informative only if $s<d(1-2\eps)$: the revealed values may induce fewer than
$d$ class changes in all. The count agrees with the parent paper's charge for a state of $q$ symbols carried across a
regime change, of order $qd$ when each of the $d$ output rows carries its own state
\cite[Section~7.2]{SeventeenthPaper2026}.

For machines the step remains open. \Cref{thm:linear-advice} uses that $\X\cap\ker\alpha$ is a code, and the proof of
\cref{thm:any-advice} cancels the offsets by subtracting a matrix of the fibre, so that each fibre lies in a coset of a
code. With nonlinear decoders an offset that depends on the row is itself a break of uniformity. So the second step of
\cite[Section~7.4]{SeventeenthPaper2026} needs, for machines, a version of \cref{thm:TP} for the cone systems of real
schedules, whose decoders are compositions of the transition function and need not be linear; a form of it that
tolerates offsets depending on the row; and a bound on the class changes that a genie induces.
\end{remark}
